\documentclass{labreport}

\usepackage{microtype}
\usepackage{enumitem}
\usepackage{pifont}
\usepackage{tabularx}
\usepackage{booktabs}

% Colors
\definecolor{codebg}{RGB}{248, 249, 251}
\definecolor{codeborder}{RGB}{218, 222, 229}
\definecolor{linkcolor}{RGB}{20, 60, 120}
\definecolor{zshgreen}{RGB}{0, 160, 80}
\definecolor{zshred}{RGB}{200, 30, 30}

% Hyperref setup with subtle link colors
\hypersetup{
    colorlinks=true,
    linkcolor=linkcolor,
    filecolor=linkcolor,
    urlcolor=linkcolor,
    citecolor=linkcolor
}

% Code and terminal listings
\lstdefinestyle{bash}{
    basicstyle=\ttfamily\footnotesize,
    backgroundcolor=\color{codebg},
    frame=single,
    framerule=0.5pt,
    rulecolor=\color{codeborder},
    breaklines=true,
    breakatwhitespace=false,
    showstringspaces=false,
    keepspaces=true,
    columns=flexible,
    numbers=none,
    xleftmargin=3pt,
    xrightmargin=3pt,
    framesep=4pt
}

\lstdefinestyle{c}{
    language=C,
    basicstyle=\ttfamily\footnotesize,
    backgroundcolor=\color{codebg},
    frame=single,
    framerule=0.5pt,
    rulecolor=\color{codeborder},
    keywordstyle=\color{blue!85!black}\bfseries,
    commentstyle=\color{gray!80!black}\itshape,
    stringstyle=\color{teal!80!black},
    numbers=left,
    numberstyle=\tiny\color{gray!80!black},
    numbersep=8pt,
    breaklines=true,
    breakatwhitespace=false,
    showstringspaces=false,
    keepspaces=true,
    columns=flexible,
    tabsize=2,
    xleftmargin=14pt,
    xrightmargin=3pt,
    framesep=4pt
}

\lstset{
    extendedchars=true,
    literate=
        {➜}{{\textcolor{zshgreen}{\ding{220}}}}{1}
        {├─}{|-\,}{2}
        {└─}{\`-\,}{2}
        {├}{|}{1}
        {─}{-}{1}
        {└}{\`}{1}
        {✗}{{\textcolor{zshred}{\ding{55}}}}{1}
}

% Section styling for clean document structure
\titleformat{\section}
  {\large\bfseries}{\thesection}{0.8em}{}[\vspace{1pt}\hrule height 0.5pt]
\titleformat{\subsection}
  {\normalsize\bfseries}{\thesubsection}{0.8em}{}
\titleformat{\subsubsection}
  {\small\bfseries}{\thesubsubsection}{0.8em}{}

\begin{document}

\labtitle
    {6}
    {Inside the xv6 File System}
    {st. gr. FAF-243}
    {Șoimu Ionuț}
    {asist. univ.\\[0.15cm] Reițman Patricia}

\tableofcontents
\newpage

% ==============================================================================
\section{Lab 6: Track B -- Inside the xv6 File System}

\subsection{Before you start}

\subsubsection*{Purpose and Scope}
The file system of the educational operating system xv6 implements the core design patterns of modern Unix storage architectures (such as Linux ext4): on-disk inode structures, directory tables mapping names to inode numbers, reference-counted hard links, indirect block addressing, and crash-recovery logging, all within a concise, readable implementation of a few hundred lines of C.

The primary objective of this laboratory work is to dissect the internal representation of files and directories in xv6. We decode the low-level attributes exposed by xv6 utilities, develop a customized stat utility (\texttt{user/mystat.c}) to track inode metadata and hard-link lifecycles, calculate theoretical and architectural filesystem storage limits from definitions in \texttt{kernel/fs.h}, and author a recursive directory traversal tool (\texttt{user/find.c}) adhering to POSIX filesystem semantics.

\subsubsection*{Environment Setup and Build Readiness}
We verify that the RISC-V cross-compilation toolchain and the QEMU emulator environment are operational, and ensure that the xv6 disk image builds cleanly:

\begin{lstlisting}[style=bash]
➜  ~ cd ~/xv6-riscv
➜  xv6-riscv git:(riscv) which riscv64-linux-gnu-gcc qemu-system-riscv64
/usr/bin/riscv64-linux-gnu-gcc
/usr/bin/qemu-system-riscv64
➜  xv6-riscv git:(riscv) make kernel/kernel fs.img
make: 'kernel/kernel' is up to date.
make: 'fs.img' is up to date.
\end{lstlisting}

% ------------------------------------------------------------------------------
\subsection{Part 1. Read what xv6 shows you}

\subsubsection{Inspecting Directory Entries and Inode Properties}
Inside xv6, running \texttt{ls} produces a four-column output displaying the entry's filename, file type identifier, inode number, and size in bytes:

\begin{lstlisting}[style=bash]
$ ls
.              1 1 1024
..             1 1 1024
README         2 2 2441
cat            2 3 35864
echo           2 4 34760
sleep          2 5 34440
args           2 6 34424
forktest       2 7 17704
grep           2 8 43288
init           2 9 35112
kill           2 10 34664
ln             2 11 34456
ls             2 12 41944
mkdir          2 13 34712
rm             2 14 34696
sh             2 15 57072
stressfs       2 16 35536
usertests      2 17 201520
grind          2 18 50336
wc             2 19 36832
zombie         2 20 34008
logstress      2 21 36680
forphan        2 22 35472
dorphan        2 23 34936
sync           2 24 34112
console        3 25 0
$ mkdir a
$ ls a
.              1 28 32
..             1 1 1024
\end{lstlisting}

\subsubsection*{Observe:}
\begin{enumerate}[leftmargin=*]
    \item \textbf{What do types 1, 2 and 3 mean (kernel/stat.h)?}\\[0.2em]
    In \texttt{kernel/stat.h}, file types are defined as:
    \begin{itemize}[noitemsep, topsep=2pt]
        \item \textbf{\texttt{T\_DIR}} (\textbf{1}): Directory -- A file containing a list of \texttt{struct dirent} entries mapping names to inode numbers.
        \item \textbf{\texttt{T\_FILE}} (\textbf{2}): Regular file -- An unformatted linear array of data bytes.
        \item \textbf{\texttt{T\_DEVICE}} (\textbf{3}): Device -- A special character or block device interface (e.g., \texttt{console}), identified by major and minor device numbers.
    \end{itemize}

    \item \textbf{Which inode number does the root directory have, and why are . and .. the same there?}\\[0.2em]
    The root directory has inode number \textbf{1} (\texttt{ROOTINO 1} in \texttt{kernel/fs.h}).\\
    In the root directory, both \texttt{.} (the current directory) and \texttt{..} (the parent directory) point to inode 1 because the root directory is the top-level anchor of the entire filesystem hierarchy; it has no parent directory above it, so by convention its parent pointer \texttt{..} points back to itself.

    \item \textbf{An empty directory a has size 32. A directory entry has a fixed size. How big is one entry, and which two entries does a contain?}\\[0.2em]
    In xv6, directory entries are defined by \texttt{struct dirent} in \texttt{kernel/fs.h}:
    \begin{lstlisting}[style=c]
    struct dirent {
      ushort inum;
      char name[DIRSIZ];
    };
    \end{lstlisting}
    Since \texttt{ushort} is 2 bytes and \texttt{DIRSIZ} is 14 bytes, each directory entry occupies exactly:
    \[
    2 \text{ bytes} + 14 \text{ bytes} = \mathbf{16 \text{ bytes}}
    \]
    A newly created directory contains exactly two mandatory self-referential entries:
    \begin{enumerate}[label=(\alph*), noitemsep]
        \item \texttt{.} (pointing to directory \texttt{a}'s own inode number, 28)
        \item \texttt{..} (pointing to parent root directory inode 1)
    \end{enumerate}
    Consequently, its total initial size is:
    \[
    2 \text{ entries} \times 16 \text{ bytes/entry} = \mathbf{32 \text{ bytes}}
    \]
\end{enumerate}

% ------------------------------------------------------------------------------
\subsection{Part 2. Write mystat}

\subsubsection{Implementation of user/mystat.c}
We implement \texttt{user/mystat.c} to invoke the \texttt{stat()} system call on each path specified as a command-line argument and print its formatted attributes:

\begin{lstlisting}[style=c]
#include "kernel/types.h"
#include "kernel/stat.h"
#include "user/user.h"

char *
typename(int t)
{
  if(t == T_DIR)
    return "directory";
  if(t == T_FILE)
    return "file";
  if(t == T_DEVICE)
    return "device";
  return "?";
}

int
main(int argc, char *argv[])
{
  struct stat st;
  if(argc < 2){
    fprintf(2, "usage: mystat <path>...\n");
    exit(1);
  }
  for(int i = 1; i < argc; i++){
    if(stat(argv[i], &st) < 0){
      fprintf(2, "mystat: cannot stat %s\n", argv[i]);
      continue;
    }
    printf("%s: %s, inode %d, links %d, size %d bytes\n",
           argv[i], typename(st.type), st.ino, st.nlink, (int)st.size);
  }
  exit(0);
}
\end{lstlisting}
We register \texttt{\$U/\_mystat\textbackslash} in \texttt{UPROGS} in the xv6 \texttt{Makefile}.

\subsubsection{Tracking Hard Links at Runtime}
We build and run \texttt{mystat} inside xv6 to inspect different file types and observe how hard links affect the underlying inode link count:

\begin{lstlisting}[style=bash]
$ mystat README . console
README: file, inode 2, links 1, size 2441 bytes
.: directory, inode 1, links 2, size 1024 bytes
console: device, inode 27, links 1, size 0 bytes

$ ln README r2
$ mystat README r2
README: file, inode 2, links 2, size 2441 bytes
r2: file, inode 2, links 2, size 2441 bytes

$ rm README
$ mystat r2
r2: file, inode 2, links 1, size 2441 bytes

$ cat r2 | wc
48 336 2441
\end{lstlisting}

\subsubsection*{Observe:}
\begin{enumerate}[leftmargin=*]
    \item \textbf{Paste the outputs after ln and after rm. What changed and what did not?}
    \begin{itemize}[noitemsep, topsep=2pt]
        \item \textbf{Output after \texttt{ln README r2}:}
\begin{lstlisting}[style=bash]
README: file, inode 2, links 2, size 2441 bytes
r2: file, inode 2, links 2, size 2441 bytes
\end{lstlisting}
        Both names share the exact same inode number (\textbf{inode 2}) and size (\textbf{2441 bytes}). The hard link count (\texttt{links}) for inode 2 incremented from 1 to \textbf{2}.
        \item \textbf{Output after \texttt{rm README}:}
\begin{lstlisting}[style=bash]
r2: file, inode 2, links 1, size 2441 bytes
\end{lstlisting}
        The name \texttt{README} was erased from the root directory table, but the inode number (inode 2) and file content/size (2441 bytes) did not change. The link count decremented back from 2 to \textbf{1}. Reading \texttt{r2} via \texttt{cat r2 | wc} returns the full 2441 bytes unaltered.
    \end{itemize}

    \item \textbf{Find the line in kernel/sysfile.c (sys\_unlink) that decrements the link count. When are the data blocks freed?}\\[0.2em]
    In \texttt{kernel/sysfile.c}, inside \texttt{sys\_unlink()}, line 244 decrements the link count:
\begin{lstlisting}[style=c]
  ip->nlink--;
  iupdate(ip);
\end{lstlisting}
    \textbf{When data blocks are freed:} Data blocks are not freed during \texttt{sys\_unlink()} if other directory references or open file handles remain. Instead, freeing occurs inside \texttt{iput()} in \texttt{kernel/fs.c}. When an inode's in-memory reference count drops to 1 (\texttt{ip->ref == 1}, meaning no running processes hold an open file descriptor to it) and its link count is zero (\texttt{ip->nlink == 0}), \texttt{iput()} calls \texttt{itrunc(ip)}:
\begin{lstlisting}[style=c]
  if(ip->valid && ip->nlink == 0){
    ...
    itrunc(ip);
    ip->type = 0;
    iupdate(ip);
  }
\end{lstlisting}
    \texttt{itrunc()} frees all 12 direct data blocks and all blocks referenced by the indirect block via \texttt{bfree()}, and zeroes out the inode.

    \item \textbf{xv6 has no symbolic links and no permissions. Look at struct dinode in kernel/fs.h: which fields does an xv6 inode have?}\\[0.2em]
    In \texttt{kernel/fs.h}, \texttt{struct dinode} represents the on-disk inode structure:
\begin{lstlisting}[style=c]
struct dinode {
  short type;              // File type (T_DIR, T_FILE, T_DEVICE)
  short major;             // Major device number (T_DEVICE only)
  short minor;             // Minor device number (T_DEVICE only)
  short nlink;             // Number of links to inode in file system
  uint size;               // Size of file (bytes)
  uint addrs[NDIRECT + 1]; // Data block addresses (12 direct + 1 indirect)
};
\end{lstlisting}
    It contains only 6 fields: \texttt{type}, \texttt{major}, \texttt{minor}, \texttt{nlink}, \texttt{size}, and \texttt{addrs}.\\
    Noticeably absent are:
    \begin{itemize}[noitemsep]
        \item No ownership fields (\texttt{uid}, \texttt{gid})
        \item No permission bitmasks (\texttt{rwxrwxrwx}, \texttt{mode})
        \item No access/modification timestamps (\texttt{atime}, \texttt{mtime}, \texttt{ctime})
        \item No symbolic link target buffer or symlink type flag
    \end{itemize}
\end{enumerate}

% ------------------------------------------------------------------------------
\subsection{Part 3. Measure the limits}

\subsubsection{Inspecting Filesystem Constants and Disk Geometry}
We inspect the structural constants in \texttt{kernel/fs.h} and examine the disk layout printed by \texttt{mkfs}:

\begin{lstlisting}[style=bash]
➜  xv6-riscv git:(riscv) grep -n "define BSIZE\|define NDIRECT\|define NINDIRECT\|define DIRSIZ" kernel/fs.h
5:#define BSIZE   1024 // block size
26:#define NDIRECT   12
27:#define NINDIRECT (BSIZE / sizeof(uint))
54:#define DIRSIZ 14

➜  xv6-riscv git:(riscv) make clean >/dev/null && make fs.img 2>&1 | grep nmeta
nmeta 47 (boot, super, log blocks 31, inode blocks 13, bitmap blocks 1) blocks 1953 total 2000
\end{lstlisting}

\subsubsection{Filename Length Boundary Test}
We test filename truncation in xv6 by creating a directory with a name exceeding 14 characters:

\begin{lstlisting}[style=bash]
$ mkdir averyveryverylongname
$ ls
.              1 1 1024
..             1 1 1024
averyveryveryl 1 29 32
cat            2 3 35864
...
\end{lstlisting}

\subsubsection*{Observe:}
\begin{enumerate}[leftmargin=*]
    \item \textbf{What does ls show for averyveryverylongname, and which constant explains it?}\\[0.2em]
    \texttt{ls} displays \textbf{\texttt{averyveryveryl}}.\\
    The string \texttt{averyveryverylongname} was truncated to exactly 14 characters. This limit is governed by the constant:
    \begin{center}
    \texttt{\#define DIRSIZ 14} \quad in \texttt{kernel/fs.h}
    \end{center}
    Because each \texttt{struct dirent} allocates 14 bytes for \texttt{name[DIRSIZ]}, characters beyond the 14th are dropped by \texttt{dirlink()}.

    \item \textbf{Compute the largest file xv6 can hold: 12 direct blocks plus one indirect block that holds block numbers of 4 bytes each, blocks of BSIZE bytes. Show the calculation.}
    \begin{itemize}[noitemsep, topsep=2pt]
        \item Block size: $BSIZE = \mathbf{1024 \text{ bytes}}$.
        \item Direct blocks: $NDIRECT = \mathbf{12 \text{ blocks}}$.
        \item Size of a block number address: $\text{sizeof}(\text{uint}) = \mathbf{4 \text{ bytes}}$.
        \item Capacity of one indirect block:
        \[
        NINDIRECT = \frac{BSIZE}{\text{sizeof}(\text{uint})} = \frac{1024}{4} = \mathbf{256 \text{ data blocks}}
        \]
        \item Total maximum block addresses per inode:
        \[
        MAXFILE = NDIRECT + NINDIRECT = 12 + 256 = \mathbf{268 \text{ blocks}}
        \]
        \item Maximum theoretical file size:
        \[
        \text{Max File Size} = 268 \text{ blocks} \times 1024 \text{ bytes/block} = \mathbf{274,432 \text{ bytes}} \approx \mathbf{268 \text{ KB}}
        \]
    \end{itemize}

    \item \textbf{From the nmeta line: how many blocks does the whole disk have, and how many are for the log, the inodes, and the data?}
    \begin{itemize}[noitemsep, topsep=2pt]
        \item \textbf{Total disk image size:} \textbf{2000 blocks} (2,048,000 bytes = 2 MB).
        \item \textbf{Metadata blocks (\texttt{nmeta 47}):}
        \begin{itemize}[noitemsep]
            \item Boot block: \textbf{1 block} (sector 0)
            \item Super block: \textbf{1 block} (sector 1)
            \item Log blocks: \textbf{31 blocks} (sectors 2--32)
            \item Inode blocks: \textbf{13 blocks} (sectors 33--45; holds $13 \times 16 = 208$ inodes)
            \item Free bitmap: \textbf{1 block} (sector 46; tracks up to $1024 \times 8 = 8192$ blocks)
        \end{itemize}
        \item \textbf{Data blocks:} $2000 - 47 = \mathbf{1953 \text{ data blocks}}$ (sectors 47--1999).
    \end{itemize}
\end{enumerate}

% ------------------------------------------------------------------------------
\subsection{Part 4. Write find}

\subsubsection{Implementation of user/find.c}
We implement \texttt{user/find.c} to recursively walk the directory tree and print any path matching a target basename. It skips unallocated entries (\texttt{de.inum == 0}), prevents infinite recursion by ignoring \texttt{.} and \texttt{..}, closes opened directory descriptors to prevent descriptor exhaustion, and supports an optional \texttt{-n} switch for displaying inode numbers:

\begin{lstlisting}[style=c]
#include "kernel/types.h"
#include "kernel/stat.h"
#include "user/user.h"
#include "kernel/fs.h"
#include "kernel/fcntl.h"

void
find(char *path, char *target, int show_inode)
{
  char buf[512], *p;
  int fd;
  struct dirent de;
  struct stat st;

  if((fd = open(path, O_RDONLY)) < 0){
    fprintf(2, "find: cannot open %s\n", path);
    return;
  }

  if(fstat(fd, &st) < 0){
    fprintf(2, "find: cannot stat %s\n", path);
    close(fd);
    return;
  }

  if(st.type != T_DIR){
    close(fd);
    return;
  }

  if(strlen(path) + 1 + DIRSIZ + 1 > sizeof(buf)){
    fprintf(2, "find: path too long\n");
    close(fd);
    return;
  }

  strcpy(buf, path);
  p = buf + strlen(buf);
  *p++ = '/';

  while(read(fd, &de, sizeof(de)) == sizeof(de)){
    if(de.inum == 0)
      continue;

    char name[DIRSIZ + 1];
    memmove(name, de.name, DIRSIZ);
    name[DIRSIZ] = 0;

    // Skip . and .. to prevent infinite recursion
    if(strcmp(name, ".") == 0 || strcmp(name, "..") == 0)
      continue;

    memmove(p, de.name, DIRSIZ);
    p[DIRSIZ] = 0;

    // Trim trailing spaces if any
    int len = strlen(p);
    while(len > 0 && p[len - 1] == ' ') {
      p[len - 1] = '\0';
      len--;
    }

    if(stat(buf, &st) < 0){
      fprintf(2, "find: cannot stat %s\n", buf);
      continue;
    }

    if(strcmp(name, target) == 0){
      if(show_inode)
        printf("%s %d\n", buf, st.ino);
      else
        printf("%s\n", buf);
    }

    if(st.type == T_DIR){
      find(buf, target, show_inode);
    }
  }

  close(fd);
}

int
main(int argc, char *argv[])
{
  int show_inode = 0;

  if(argc < 3){
    fprintf(2, "usage: find <path> <name> [-n]\n");
    exit(1);
  }

  if(argc >= 4 && strcmp(argv[3], "-n") == 0)
    show_inode = 1;

  find(argv[1], argv[2], show_inode);
  exit(0);
}
\end{lstlisting}
We register \texttt{\$U/\_find\textbackslash} in \texttt{UPROGS} in \texttt{Makefile}.

\subsubsection{Validation on Nested Directories and Hard Links}
We boot xv6 on a clean disk and execute the standard test sequence:

\begin{lstlisting}[style=bash]
$ mkdir a
$ mkdir a/b
$ echo x > a/b/r2
$ echo y > a/r2
$ ln cat r2

$ find . r2
./a/b/r2
./a/r2
./r2

$ find a b
a/b
\end{lstlisting}

\subsubsection*{Stretch Goal: Inode Reporting with -n}
We execute \texttt{find} with the \texttt{-n} flag to display the inode numbers alongside each discovered path:

\begin{lstlisting}[style=bash]
$ find . r2 -n
./a/b/r2 31
./a/r2 32
./r2 3
\end{lstlisting}

\subsubsection*{Observe:}
\textbf{Which of the three r2 files are the same file as cat?}\\[0.2em]
The file \textbf{\texttt{./r2}} (with inode number \textbf{3}) is the exact same underlying file as \texttt{cat} (which also has inode number 3). Although all three files share the same leaf filename (\texttt{r2}), \texttt{./a/b/r2} (inode 31) and \texttt{./a/r2} (inode 32) are distinct, newly created regular files with unique inode entries and independent disk blocks. In contrast, \texttt{./r2} was created via \texttt{ln cat r2}, which created a second directory entry in \texttt{.} pointing to \texttt{cat}'s existing inode 3.

% ==============================================================================
\section{Conclusion}

In this laboratory work, I investigated the foundational mechanics of the Unix file system through the implementation of xv6.

By reading the on-disk structures and decoding the output of \texttt{ls}, I explored how a hierarchical directory tree is built out of flat tables of \texttt{struct dirent} structures that map human-readable names to numeric inode references. Developing \texttt{mystat} made concrete how hard links function: creating a link merely registers an additional name pointing to an existing inode and increments its \texttt{nlink} counter, while unlinking removes the directory entry and decrements the counter without freeing data blocks until all active directory links and open file descriptors are released.

Analyzing \texttt{kernel/fs.h} illuminated the structural trade-offs of educational file system design: fixed 14-character filename boundaries (\texttt{DIRSIZ}), a 1024-byte block size (\texttt{BSIZE}), and a 12-direct plus 1-indirect block addressing scheme that limits individual file sizes to approximately 268 KB. Finally, implementing \texttt{find} demonstrated the practical mechanics of directory traversal, emphasizing the critical importance of skipping recursive loops (\texttt{.} and \texttt{..}) and managing process file descriptor limits. Together, these exercises demystified the operating system storage stack from high-level user utilities down to raw disk sectors.

\end{document}
